\scanpage{71}\mspage{58}
% archive sheet number: 70

b) $C$ et $TC=C\cap W$ sont stables par composition
transfinie.

Supposons de plus que $u^*$ applique $M(I'^{\unsure{o|\circ}})_c$
dans $M(I^o)_c$, ce qui est le cas notamment
si $u$ est loc. \deleted{\unsure{filtr.}} cofiltrant i.e. les
\deleted{\unsure{cat.}}\inserted{ordonnés} localisés $\mathord{}_{i'}\backslash I$ sont loc. cofiltrants. \deleted{\unsure{à g.}}
Alors le foncteur $u_!^W$ construit en (\unsure{3.3.1|3.2.1}).
est adjoint à g. de $u_W^*$.

\uline{Remarque} 3.6. Quand on a un couple
de foncteurs adjoints

\diagram{\begin{tikzcd}[ampersand replacement=\&,column sep=large]
C \arrow[r,dash,shift left=1ex,"u_!", "\unsure{\text{---}|\to}"{description,near end,above=1pt,font=\tiny}] \& C', \arrow[l,shift left=1ex,"u^*"]
\end{tikzcd}}

avec des localiseurs $W,W'$ de $C,C'$,
$u^*$ compatible aux localiseurs, et si de
plus on a une sous-catégorie \inserted{pleine} $C_0$ de $C$
telle que a) $\deleted{\unsure{C}}u_!|C_0$ transforme flèches
de $W$ en flèches de $W'$ et b) l'inclusion
$C_0\hookrightarrow C$ induit une \uline{équivalence}
$C_0W_0^{-1}\overset{\approx}{\text{---}}CW^{-1}$ (où $W_0=W\cap\mathrm{Fl}(C_0)$),
on trouve deux foncteurs canoniques

\diagram{\begin{tikzcd}[ampersand replacement=\&,column sep=large]
C W^{-1} \arrow[r,dash,shift left=1.2ex,"u_!^W"] \& C' W'^{-1}, \arrow[l,shift left=1.2ex,"u_W^*"]
\end{tikzcd}}

où $u_!^W$ est défini \inserted{seulement} à isom \unsure{can$^{\mathrm{l}}$|can$'$} près.
